% The second elementary route: blindness, and the two-site positive eigenvector.
% Lane: O (operator bridge), modules YangMills/OS/**.
% REGIME: tricotomy; NO scores, NO desk language, NO commit chronology beyond
% the reproducibility freeze.
\documentclass[11pt]{article}
\usepackage[margin=1.1in]{geometry}
\usepackage{amsmath,amssymb,amsthm}
\usepackage{booktabs,longtable,array}
\usepackage[colorlinks=true,linkcolor=blue,urlcolor=blue,citecolor=blue]{hyperref}

\newtheorem{theorem}{Theorem}[section]
\newtheorem{lemma}[theorem]{Lemma}
\newtheorem{proposition}[theorem]{Proposition}
\newtheorem{corollary}[theorem]{Corollary}
\newtheorem{remark}[theorem]{Remark}
\theoremstyle{definition}
\newtheorem{definition}[theorem]{Definition}

\emergencystretch=3em
\tolerance=2000

\newcommand{\lean}[1]{\texttt{#1}}
\newcommand{\anchor}{f181ed90e93f432159a46afd703c8f3b9898f6d2}
\newcommand{\osline}[3]{%
  \href{https://github.com/lluiseriksson/THE-ERIKSSON-PROGRAMME/blob/\anchor/YangMills/OS/#1.lean\#L#2}{\nolinkurl{#3}}}
\newcommand{\repofile}[2]{%
  \href{https://github.com/lluiseriksson/THE-ERIKSSON-PROGRAMME/blob/\anchor/#1}{\nolinkurl{#2}}}
\newcommand{\R}{\mathbb{R}}
\newcommand{\Z}{\mathbb{Z}}
\newcommand{\SU}{\mathrm{SU}}
\newcommand{\Om}{\Omega}

\title{Blind to the Coupling:\\
a Second Machine-Checked Obstruction at Spatial Extent\\
\large and the strictly positive eigenvector the elementary route failed to
produce, in closed form at two sites}
\author{Lluis Eriksson}
\date{}

\begin{document}
\maketitle

\begin{abstract}
A companion paper proved that when a spatial coupling is switched on in a
$\Z_2$ lattice gauge slice, the transfer kernel loses constant row sums, so the
uniform vector is no longer fixed and the elementary route to the vacuum
stops.  The standard replacement, when row sums fail, is the Hilbert projective
metric: a strictly positive kernel contracts it, and the contraction factor
bounds the subdominant spectral ratio.  This paper asks what that replacement
gives here, and answers in Lean~4 with \texttt{mathlib}.

It gives nothing, for two independent reasons, and both are proved.  First,
\emph{blindness}: the projective cross-ratio is invariant under multiplication
by any nowhere-zero function of the source configuration alone.  The coupled
kernel is exactly such a product, so the metric assigns the interacting and the
non-interacting kernels the same diameter at every spatial extent --- the route
cannot see the coupling at all.  Second, \emph{volume degeneration}: two
constant configurations realise the cross-ratio $e^{4\beta L}$, so every
admissible projective diameter is at least $4\beta L$ and every contraction
factor obtainable this way is at least $\tanh(\beta L)$, which lies within
$2e^{-2\beta L}$ of the trivial bound $1$.  At the one place where the truth is
known --- the decoupled kernel, whose subdominant ratio the companion paper
computes to be exactly $\tanh\beta$ at every $L$ --- this route already returns
$\tanh(\beta L)$ instead.  The degeneration is the method's, not the model's.

We then hand over the object the elementary route stopped producing, at the
smallest interacting size.  In the character basis the coupled two-site kernel
splits into two $2\times2$ blocks, and we exhibit a \emph{strictly positive}
eigenvector in closed form, together with a second exact eigenpair.  The
identity $A-B=4$ between the two even-sector eigenvalues drives every estimate.

\textbf{No volume-uniform statement about an interacting system is proved here,
and none is claimed}; the general-$L$ behaviour is recorded separately as
measured and unproved.  Nothing in this paper is a claim about $\SU(N)$, the
continuum limit, or the Yang--Mills mass gap.
\end{abstract}

\section{Introduction}

\subsection{What the predecessor left}

The companion paper \cite{ErikssonSpatial} proved, for a $\Z_2$ slice carrying
$L$ spins, that the decoupled transfer kernel has constant row sums at every
$L$ --- so the uniform vacuum survives --- and that switching on a coupling
inside the slice destroys that constancy, because the spatial weight depends on
the \emph{source} only and therefore factors out of the sum over the target.
With constant row sums goes $T\Om=\Om$, and with it the starting point of every
later step of the reconstruction.

That paper ended by naming two honest continuations: a Perron--Frobenius route
to the vacuum at spatial extent, or a statement that the elementary methods end
there.  It declined to say which.  This paper answers, and the answer has both
signs: the natural elementary replacement is proved useless here, and the
object it failed to produce is written down at the smallest interacting size ---
as a strictly positive eigenvector in closed form, which is what is proved about
it.

\subsection{The natural next attempt, and why it is the right one to test}

When constant row sums fail, the textbook instrument is the \emph{Hilbert
projective metric}.  A strictly positive kernel contracts it, the contraction
factor is $\tanh(\Delta/4)$ where $\Delta$ is the projective diameter
\cite{Birkhoff,Bushell}, and that factor bounds the subdominant spectral ratio
without any need to know the leading eigenvector.  It is exactly the tool one
reaches for when the eigenvector is unavailable, which is the situation
\cite{ErikssonSpatial} created.  So it is the right next thing to test, and
testing it is what this paper does.

\subsection{What is proved}

Write $k_\beta$ for the Boltzmann bond weight, $Z_\beta = e^\beta + e^{-\beta}$,
$K(\sigma,\tau)=\prod_j k_\beta(\sigma_j,\tau_j)$ for the decoupled kernel, and
$w$ for a weight depending on the source configuration only.  Define the
projective cross-ratio
\[
  R_A(\sigma,\sigma';\tau,\tau') \;=\;
  \frac{A(\sigma,\tau)\,A(\sigma',\tau')}{A(\sigma,\tau')\,A(\sigma',\tau)} .
\]

\medskip
\noindent\textbf{Blindness.}
\begin{itemize}
\item $R_{wK} = R_K$ for \emph{every} nowhere-zero source-only weight $w$, at
      every spatial extent.  The weight occurs once above and once below the
      bar, at the same two configurations, and cancels identically.
\end{itemize}

\medskip
\noindent\textbf{Volume degeneration.}
\begin{itemize}
\item at the all-$0$ and all-$1$ configurations, $R_K = e^{4\beta L}$ exactly;
\item hence every $D$ with $\log R_K \le D$ everywhere satisfies $4\beta L \le
      D$, and every contraction factor $\tanh(D/4)$ obtainable this way
      satisfies $\tanh(\beta L)\le \tanh(D/4)$;
\item and $1-\tanh(\beta L)\le 2e^{-2\beta L}$: that lower bound is itself
      exponentially close, in the volume, to the trivial bound $1$;
\item while the truth for the decoupled kernel is $\tanh\beta$ at every $L$
      \cite{ErikssonSpatial}, strictly better for $L\ge2$, $\beta>0$.
\end{itemize}

\medskip
\noindent\textbf{The two-site positive eigenvector.}  With
$A=(e^\beta+e^{-\beta})^2$ and $B=(e^\beta-e^{-\beta})^2$:
\begin{itemize}
\item the coupled kernel preserves the two-dimensional span of the constant and
      the top character, acting there by
      $\bigl(\begin{smallmatrix} A\cosh\gamma & B\sinh\gamma\\
      A\sinh\gamma & B\cosh\gamma\end{smallmatrix}\bigr)$;
\item the vector $\alpha\cdot 1+\delta\cdot s_0s_1$ with
      $\alpha = B\sinh\gamma$ and
      $\delta=\lambda_+-A\cosh\gamma$ is an eigenvector with eigenvalue
      $\lambda_+ = \tfrac12\bigl((A+B)\cosh\gamma+\sqrt{(A+B)^2\cosh^2\gamma-4AB}\bigr)$,
      and it is \emph{strictly positive} at every configuration for
      $\beta>0$, $\gamma>0$;
\item $s_0+s_1$ is a second eigenvector, with eigenvalue
      $\mu_+ = (e^\beta-e^{-\beta})Z_\beta e^{\gamma}$, and it does not vanish;
\item $\mu_+<\lambda_+$ for $\beta>0$, $\gamma>0$, so the two exhibited
      eigenpairs are ordered.  The identity $A+B-2(e^\beta-e^{-\beta})Z_\beta
      =4e^{-2\beta}$, companion of $A-B=4$, reduces that to one line.
\end{itemize}

\subsection{What is deliberately not claimed}

\begin{enumerate}
\item \textbf{No volume-uniform statement about an interacting system.}  None is
proved here and none is claimed.  Section~\ref{sec:measured} records what is
\emph{measured} about general $L$ and says explicitly that it is not proved.
\item \textbf{No claim that contraction arguments cannot work.}  What is proved
is that this one --- the Hilbert metric applied to this kernel --- is blind and
degenerate.  A different metric, a power of the kernel, or a Dobrushin-type
condition are untouched by these theorems.
\item \textbf{The two-site eigenvector is not claimed to be the spectral
radius.}  It is a strictly positive eigenvector.  That a strictly positive
eigenvector of a positive matrix carries the spectral radius is the
Perron--Frobenius theorem; that classical identification is \emph{not}
machine-checked in this development, and Remark~\ref{rem:pf} says so.
\item \textbf{No completeness of the two-site spectrum.}  Two eigenpairs are
exhibited.  That they and their partners exhaust the four-dimensional space is
not formalised.
\item \textbf{Still $\Z_2$}, one variable per site, and the closed form is at
$L=2$.
\item \textbf{Nothing about $\SU(N)$, the continuum limit, or the Clay problem}
\cite{Clay}.  No reduction is claimed and none is known to the author.
\end{enumerate}

\section{The projective cross-ratio}

\begin{definition}\label{def:cr}
For a kernel $A:X\times X\to\R$ on a set $X$,
\[
  R_A(\sigma,\sigma';\tau,\tau') \;=\;
  \frac{A(\sigma,\tau)\,A(\sigma',\tau')}{A(\sigma,\tau')\,A(\sigma',\tau)} .
\]
\end{definition}

\osline{SpatialBirkhoff}{121}{crossRatio}.  The Hilbert projective diameter of
a strictly positive kernel is $\sup\log R_A$, and the Birkhoff contraction
factor is $\tanh(\Delta/4)$ \cite{Birkhoff,Bushell}.  Everything below is
phrased so that only $R_A$ itself is formalised: the two classical facts just
quoted are used as \emph{context}, never as hypotheses of a theorem.

\section{Blindness: no source-only weight is visible}

\begin{theorem}[The metric cannot see a source-only weight]\label{thm:blind}
Let $w$ be nowhere zero and depend on the source only, and let
$(wK)(\sigma,\tau) = w(\sigma)K(\sigma,\tau)$.  Then
$R_{wK} = R_{K}$ at every four configurations, at every spatial extent.
\end{theorem}

\begin{proof}
$w(\sigma)$ and $w(\sigma')$ each occur once in the numerator and once in the
denominator, so they cancel.
\end{proof}

\osline{SpatialBirkhoff}{136}{crossRatio\_sourceWeighted}, and for the
coupled kernel of \cite{ErikssonSpatial} as the special case
\osline{SpatialBirkhoff}{156}{crossRatio\_coupledKernel}.

\begin{remark}
The hypothesis is not a technicality chosen to make the proof work: it is
\emph{exactly} the structural feature that \cite{ErikssonSpatial} used to break
the row sums.  The same one-line fact --- the weight depends on the source only
--- destroys the first route and makes the second route unable to detect that
anything happened.  One property, two opposite-looking failures.
\end{remark}

\section{Volume degeneration}

\begin{theorem}[The witness]\label{thm:witness}
$R_K$ at the all-$0$ and all-$1$ configurations equals $e^{4\beta L}$.
\end{theorem}

\begin{proof}
The four kernel values are $e^{\beta L}$, $e^{\beta L}$, $e^{-\beta L}$,
$e^{-\beta L}$.
\end{proof}

\osline{SpatialBirkhoff}{192}{crossRatio\_const}.

\begin{theorem}[The wall]\label{thm:wall}
If $D$ satisfies $\log R_K \le D$ at every four configurations, then
$4\beta L\le D$, and consequently $\tanh(\beta L)\le\tanh(D/4)$.
\end{theorem}

\osline{SpatialBirkhoff}{207}{projDiameter\_ge},
\osline{SpatialBirkhoff}{216}{birkhoff\_bound\_ge}.

\begin{theorem}[The degeneration is exponential in the volume]\label{thm:degen}
For $\beta\ge0$, $\;1-\tanh(\beta L)\le 2e^{-2\beta L}$.
\end{theorem}

\osline{SpatialBirkhoff}{228}{birkhoff\_bound\_near\_one}.

\begin{theorem}[And the bound is not tight]\label{thm:nottight}
For $\beta>0$ and $L\ge2$, $\;\tanh\beta<\tanh(\beta L)$.
\end{theorem}

\osline{SpatialBirkhoff}{237}{birkhoff\_bound\_not\_tight}.

\begin{remark}
Theorem~\ref{thm:nottight} is what makes Theorem~\ref{thm:degen} a statement
about the \emph{method}.  For the decoupled kernel the companion paper computes
the subdominant ratio exactly --- it is $\tanh\beta$, independent of $L$ --- so
in the one case where the answer is known, the route already returns a number
that degenerates while the truth does not.  Nothing about the model forces the
loss; the instrument loses it.
\end{remark}

\section{The positive eigenvector, in closed form at two sites}

The obstruction of \cite{ErikssonSpatial} says the vacuum stops being handed
to us.  It does not say the vacuum is unavailable.  At $L=2$ the object it
stopped handing over can be written down, and this section does so --- as what
is actually proved about it, an eigen-equation together with strict positivity;
Remark~\ref{rem:pf} says precisely what would be needed to call it the vacuum
and why that step is absent.

\begin{definition}
$s_i(\sigma)$ is the sign of site $i$, and $s_0s_1$ is the top character.  The
four characters $1,s_0,s_1,s_0s_1$ span the functions on a two-site slice.
\end{definition}

\begin{lemma}[The weight is affine in the top character]\label{lem:affine}
$w_\gamma(\sigma)=\cosh\gamma+\sinh\gamma\,(s_0s_1)(\sigma)$.
\end{lemma}

\osline{SpatialPerron}{153}{spatialWeight\_eq}.  The exponent takes only the
values $\pm\gamma$, which is the whole content.

\begin{lemma}[The decoupled kernel is diagonal in this basis]\label{lem:diag}
$\sum_\tau K(\sigma,\tau)=Z_\beta^{L}$,
$\sum_\tau K(\sigma,\tau)s_i(\tau)=(e^\beta-e^{-\beta})Z_\beta^{L-1}s_i(\sigma)$
\emph{\cite{ErikssonSpatial}}, and
$\sum_\tau K(\sigma,\tau)(s_0\cdots s_{L-1})(\tau)
 = (e^\beta-e^{-\beta})^{L}(s_0\cdots s_{L-1})(\sigma)$.
\end{lemma}

\osline{SpatialPerron}{132}{spatialKernel\_allSign} for the third; the first
two are theorems of the companion development.

\begin{theorem}[The even sector closes]\label{thm:even}
With $A=Z_\beta^2$ and $B=(e^\beta-e^{-\beta})^2$, the coupled kernel maps
$\alpha\cdot1+\delta\cdot s_0s_1$ to
$(\cosh\gamma\,\alpha A+\sinh\gamma\,\delta B)\cdot 1
 +(\cosh\gamma\,\delta B+\sinh\gamma\,\alpha A)\cdot s_0s_1$.
\end{theorem}

\osline{SpatialPerron}{194}{coupled\_even\_action}.

\begin{lemma}[The identity that drives every estimate]\label{lem:AB}
$A-B=4$.
\end{lemma}

\osline{SpatialPerron}{183}{evenA\_sub\_evenB}.  Indeed
$(a+b)^2-(a-b)^2=4ab$ and $e^{\beta}e^{-\beta}=1$.

\begin{theorem}[A strictly positive eigenvector, in closed form]\label{thm:vacuum}
Let $\lambda_+=\tfrac12\bigl((A+B)\cosh\gamma+\sqrt{(A+B)^2\cosh^2\gamma-4AB}\bigr)$.
Then $\alpha=B\sinh\gamma$, $\delta=\lambda_+-A\cosh\gamma$ gives an
eigenvector of the coupled kernel with eigenvalue $\lambda_+$, and for
$\beta>0$ and $\gamma>0$ that eigenvector is strictly positive at every
configuration.
\end{theorem}

\osline{SpatialPerron}{249}{coupled\_perron\_eigen},
\osline{SpatialPerron}{261}{perronVec\_pos}; the discriminant is nonnegative
by \osline{SpatialPerron}{216}{evenDisc\_nonneg}, using $\cosh\gamma\ge1$.

\begin{proof}[Proof of positivity]
The vector takes the two values $\alpha\pm\delta$, so positivity is
$|\,2\delta\,|<2\alpha$, i.e.
$4\cosh\gamma-2B\sinh\gamma<\sqrt{\Delta}<4\cosh\gamma+2B\sinh\gamma$ after
using $A-B=4$.  Squaring, the two gaps are
$16B\sinh\gamma(\sinh\gamma+\cosh\gamma)$ and
$16B\sinh\gamma(\cosh\gamma-\sinh\gamma)$, both positive.
\end{proof}

\begin{theorem}[A second exact eigenpair]\label{thm:odd}
$s_0+s_1$ is an eigenvector with eigenvalue
$(e^\beta-e^{-\beta})Z_\beta e^{\gamma}$, and it takes the value $2$ at the
constant configuration, so it does not vanish identically.
\end{theorem}

\osline{SpatialPerron}{310}{coupled\_odd\_eigen},
\osline{SpatialPerron}{304}{oddVec\_cfgFlat}.

\begin{theorem}[The two exhibited eigenpairs are ordered]\label{thm:order}
For $\beta>0$ and $\gamma>0$, $\;\mu_+<\lambda_+$.
\end{theorem}

\osline{SpatialPerron}{358}{oddEigen\_lt\_evenTop}.  The proof squares once
and uses the identity
$A+B-2(e^\beta-e^{-\beta})Z_\beta = 4e^{-2\beta}$
(\osline{SpatialPerron}{347}{evenSum\_sub\_two\_odd}) together with
$AB=\bigl((e^\beta-e^{-\beta})Z_\beta\bigr)^2$
(\osline{SpatialPerron}{339}{evenA\_mul\_evenB}); the resulting gap is
$4\mu_0\cosh\gamma\cdot 4e^{-2\beta}(\cosh\gamma+\sinh\gamma)>0$.  So the ratio
$\mu_+/\lambda_+$ is a genuine number strictly below $1$ --- between the two
eigenpairs this paper exhibits.  Whether it is \emph{the} subdominant ratio
depends on completeness of the spectrum, which is not formalised.

\begin{remark}[What Perron--Frobenius would add, and that we do not use it]
\label{rem:pf}
For a strictly positive matrix, Perron--Frobenius says a strictly positive
eigenvector is unique up to scale and its eigenvalue is the spectral radius ---
so Theorem~\ref{thm:vacuum} would identify $\lambda_+$ as the top of the
spectrum and the vector as \emph{the} vacuum.  We do not invoke that, and the
reason is worth recording because it is a fact about the tools rather than about
the mathematics.  The \texttt{mathlib} revision pinned by this development
carries \emph{no} Perron--Frobenius theorem.  What it carries is the layer
underneath: \texttt{Mathlib.LinearAlgebra.Matrix.Irreducible}, a quiver-theoretic
treatment of irreducibility and primitivity for entrywise-nonnegative matrices.
The statement that an irreducible nonnegative matrix has a positive eigenvector
whose eigenvalue is the spectral radius is not there, and neither is the Hilbert
projective metric nor the Birkhoff contraction coefficient --- which is why
Definition~\ref{def:cr} builds the cross-ratio from nothing.  So the
identification would have to be constructed before it could be used, and this
paper does not construct it.  What is machine-checked is exactly what is stated:
an eigen-equation, and strict positivity of the eigenvector.
\end{remark}

\section{What is measured and not proved}\label{sec:measured}

A reader will ask what happens for $L>2$.  We have looked, numerically, and we
report the result with its label attached and no theorem behind it.

Exact dense diagonalisation of the $2^L\times2^L$ coupled kernel over
$\beta\in\{0.3,0.8\}$, $\gamma\in\{0,0.4,1.2\}$, $L\in\{2,3,4\}$ gives a
subdominant ratio that rises towards $1$ with $L$ over much of that grid --- for
instance $0.908, 0.977, 0.994$ at $\beta=0.8$, $\gamma=1.2$.  The rows where it
stays bounded away from $1$ are those with $\sinh 2\beta\,\sinh 2\gamma<1$.
Both statements are \textbf{verified} in the numerical sense and \textbf{not
proved}; the identification of that threshold with the Kramers--Wannier
disordered line of the two-dimensional Ising transfer matrix is standard
literature and is not established here.  The probe and its transcript are in the
repository, and their consequence for the programme is recorded there too: a
volume-uniform statement about the coupled system, asserted without an explicit
smallness hypothesis, is aiming at something the numbers say is false.

Nothing in this section is used by any theorem of this paper.

\section{Formalisation notes}

\paragraph{One structural fact, two failures.}  That the spatial weight depends
on the source only is what broke the row sums in \cite{ErikssonSpatial} and what
makes the projective metric blind here.  The same hypothesis appears in both
proofs, with opposite-looking consequences.

\paragraph{Characters instead of enumeration.}  The two-site computation never
enumerates the four configurations.  Writing every vector in the character basis
turns each sum into an instance of the companion module's factorisation lemma,
so the proofs are the same length at $L=2$ as they would be at any $L$ where the
sector closed.

\paragraph{One identity carries the estimates.}  $A-B=4$
(Lemma~\ref{lem:AB}) reduces both positivity inequalities to
$16B\sinh\gamma(\cosh\gamma\pm\sinh\gamma)>0$.  Without it the same
inequalities are opaque to automation.

\paragraph{Cofactors computed, not guessed.}  The coefficients of the
\texttt{linear\_combination} steps and the two square-root estimates were
obtained by polynomial division in a computer algebra system before being
written into Lean.  The division also returned something useful: the ordering
identity of Theorem~\ref{thm:order} needs only $\cosh^2-\sinh^2=1$ and
\emph{not} $e^{\beta}e^{-\beta}=1$, which removed a hypothesis from the
estimate.  Guessing those coefficients by trial would have cost more compiler
round-trips than the algebra cost.

\paragraph{Self-contained hyperbolic facts, and a naming tax.}
\texttt{cosh}/\texttt{sinh} lemmas were re-derived from \texttt{Real.cosh\_eq}
and \texttt{Real.sinh\_eq} rather than searched for by name, after the search
cost more than the derivation.  Where a library lemma was unavoidable the name
still had to be read out of the pinned source: the division-comparison lemmas
this development needs carry a \texttt{\textsubscript{0}} suffix in this
revision, and one \texttt{simp} normal form silently undoes
$e^{\gamma}=\cosh\gamma+\sinh\gamma$, so that rewrite has to be carried as a
term in the linear combination instead.

\section{Reproducibility}\label{sec:repro}

All artifacts are at source checkpoint \texttt{\anchor} on branch \texttt{main}
of \url{https://github.com/lluiseriksson/THE-ERIKSSON-PROGRAMME}.  Every link
was resolved against that commit, and checked to point at an actual declaration
line, before compilation.

\begin{center}
\begin{tabular}{ll}
\toprule
Quantity & Value \\
\midrule
Full core build & 8426 jobs, success \\
Oracle commands & 2534 \\
\quad with axiom dependencies & 2511 \\
\quad axiom-free & 23 \\
Distinct axioms across all outputs & \texttt{propext}, \texttt{Quot.sound}, \texttt{Classical.choice} \\
Occurrences of \texttt{sorryAx} & 0 \\
Project axioms & 0 \\
Errors & 0 \\
\bottomrule
\end{tabular}
\end{center}

The design probe of Section~\ref{sec:measured} is
\repofile{scripts/probe\_spatial\_birkhoff.py}{scripts/probe\_spatial\_birkhoff.py}
with transcript
\repofile{docs/O-LANE-PROBE-TRANSCRIPT-20260728.txt}{docs/O-LANE-PROBE-TRANSCRIPT-20260728.txt};
the continuation ladder it fixed is
\repofile{docs/O-LANE-CONTINUATION-20260728.md}{docs/O-LANE-CONTINUATION-20260728.md}.
Verdicts are in
\repofile{docs/VERIFICATION-LEDGER.md}{docs/VERIFICATION-LEDGER.md}.  The oracle
script is \repofile{oracle\_check.lean}{oracle\_check.lean}.

\begin{thebibliography}{9}

\bibitem{ErikssonSpatial}
L.~Eriksson,
\emph{Where the Elementary Reconstruction Stops: Spatial Coupling Breaks the
Uniform Vacuum, Machine-Checked}, viXra, 2026.

\bibitem{Birkhoff}
G.~Birkhoff,
\emph{Extensions of Jentzsch's theorem},
Transactions of the American Mathematical Society \textbf{85} (1957), 219--227.

\bibitem{Bushell}
P.~J.~Bushell,
\emph{Hilbert's metric and positive contraction mappings in a Banach space},
Archive for Rational Mechanics and Analysis \textbf{52} (1973), 330--338.

\bibitem{ErikssonOSChain}
L.~Eriksson,
\emph{From the Gibbs Weight to the Spectral Gap: A Complete Machine-Checked
Osterwalder--Seiler Chain for the $\Z_2$ Lattice Gauge Chain}, viXra, 2026.

\bibitem{Lean4}
L.~de~Moura and S.~Ullrich,
\emph{The Lean 4 theorem prover and programming language},
CADE-28, Lecture Notes in Computer Science \textbf{12699} (2021), 625--635.

\bibitem{Mathlib}
The mathlib Community,
\emph{The Lean mathematical library},
CPP 2020, 367--381.

\bibitem{Clay}
A.~Jaffe and E.~Witten,
\emph{Quantum Yang--Mills theory},
Clay Mathematics Institute Millennium Prize Problem description, 2000.

\end{thebibliography}

\end{document}
